
\chapter{死锁}

\section{死锁概念}
\concept{死锁定义}
\begin{points}
  \item 死锁是多个进程因竞争系统资源而造成的一种僵局，若无外力作用，这些进程都将永远不能向前推进。
  \item 若进程集合中的每个进程都在等待只能由该集合中其他进程引发的事件，则该集合发生死锁。
  \item 典型例子：P1 占有打印机等待读卡机，P2 占有读卡机等待打印机，双方都无法继续。
\end{points}
\plain{定义题要抓“相互等待 + 无外力就永远推进不了”这两层，重点不是普通阻塞，而是整个等待环自己解不开。}
\exam{名词解释或判断题常拿“暂时阻塞”和“死锁”混淆，答案里最好点出“若无外力作用将永远不能推进”。}


\concept{死锁与饥饿、死循环}
\begin{points}
  \item 死锁：一组进程相互等待，若无外力永远不能推进。
  \item 饥饿：某进程长期得不到所需资源，但系统整体仍在推进。
  \item 死循环：一个进程自身逻辑错误导致反复执行，不一定涉及资源等待。
  \item 活锁：进程没有阻塞，但不断改变状态、相互让步，导致没有实际进展。
\end{points}
\plain{区分题的关键是看“系统整体还动不动”：饥饿是别人还能跑，死循环是自己瞎转，死锁则是一组进程彼此卡死。}
\exam{很适合出选择/判断：给一个现象，让你判断是死锁、饥饿、活锁还是死循环。}


\section{资源分类与死锁必要条件}
\concept{资源分类}
\begin{points}
  \item 可剥夺资源：可从进程手中强行取回且不破坏计算，例如 CPU、内存页框。
  \item 不可剥夺资源：一旦分配，不能强行剥夺，否则会造成错误，例如打印机、部分外设。
  \item 可重用资源：使用完可再次分配，如设备、内存。
  \item 消耗性资源：使用后消失，如消息、信号。
  \item 竞争永久性资源和消耗性资源都可能产生死锁；竞争可剥夺资源通常不会形成典型死锁。
\end{points}
\plain{这里别把死锁只和“打印机这类永久资源”绑死，PPT 明确强调消息这类消耗性资源也能形成循环等待。}
\exam{常考可剥夺/不可剥夺、永久性/消耗性资源的分类，以及哪些竞争可能引发死锁。}


\concept{死锁产生的四个必要条件}
\begin{points}
  \item \key{互斥条件}：资源一次只能被一个进程使用。
  \item \key{请求和保持条件}：进程已保持至少一个资源，同时又请求新的资源。
  \item \key{不可剥夺条件}：资源不能被强行夺走，只能由占有进程主动释放。
  \item \key{循环等待条件}：存在进程资源等待环，$P_1$ 等 $P_2$，$P_2$ 等 $P_3$，$\ldots$，$P_n$ 等 $P_1$。
\end{points}
\plain{四个条件不是“死锁特征列表”，而是判断能否从制度上预防死锁的抓手: 只要破坏其中一个，死锁就不成立。}
\exam{简答题常要求逐条解释四个必要条件，或反问某种措施到底破坏了哪一个条件。}

\concept{死锁的根本原因与诱因}
\begin{points}
  \item 根本原因是竞争资源，而资源又不能同时满足所有进程的请求。
  \item 诱因是进程推进顺序不当，即资源申请和释放次序安排得不好。
  \item 只有“竞争资源”还不一定死锁；还要配合不合适的推进顺序才会卡死。
\end{points}
\plain{资源紧张只是埋雷，申请顺序乱了才真的把雷踩爆。}
\exam{常考简答：死锁为什么不是“有竞争就必死锁”，而是资源条件和推进顺序共同作用。}

\section{处理死锁的基本方法}
\concept{四类方法总览}
\begin{points}
  \item 死锁预防：通过制度性限制破坏死锁必要条件之一。
  \item 死锁避免：运行时根据资源需求判断分配是否安全，典型算法是银行家算法。
  \item 死锁检测与解除：允许死锁发生，定期检测，发生后恢复。
  \item 鸵鸟策略：忽略死锁问题，适用于死锁概率很低且处理成本过高的系统。
\end{points}
\plain{这题最容易混的是“预防、避免、检测”三个层次：预防是事前禁规则，避免是分配前看安全性，检测是先允许出事再查。}
\exam{常考四种处理思路的区别，尤其银行家算法属于避免，不属于检测。}


\concept{死锁预防}
\begin{points}
  \item 破坏互斥条件：把独占设备改造成共享设备，例如 SPOOLing 把打印机虚拟成共享设备。但并非所有资源都能共享。
  \item 破坏请求和保持条件：进程运行前一次性申请所有资源，或申请新资源前释放已占资源。缺点是资源利用率低、可能饥饿。
  \item 破坏不可剥夺条件：资源申请失败时剥夺其已占资源。适用于 CPU、内存等可剥夺资源，不适合打印机等。
  \item 破坏循环等待条件：对资源统一编号，进程按编号递增顺序申请资源。
\end{points}
\plain{预防题不是背四条措施，而是要能明确说出“这条办法破坏了哪个必要条件，代价是什么”。}
\exam{常考静态分配法/有序资源分配法为什么能预防死锁，答案里最好点名破坏的条件。}


\concept{有序资源分配法为什么能防止死锁}
\begin{points}
  \item 给所有资源类型统一编号，要求进程只能按编号递增顺序申请资源。
  \item 这样若进程已占有编号较小的资源，就不能再回头去申请编号更小或相同层次的资源。
  \item 因而不可能形成“资源编号沿环路严格递增、最后又回到起点”的循环等待。
  \item 所以该方法本质上是通过破坏循环等待条件来预防死锁。
\end{points}
\plain{考试答“为什么有效”时不要只写一句“按顺序申请”，要把“因此不可能形成环”这层逻辑补出来。}
\exam{常考证明式简答：为什么有序资源分配法可以防止死锁。}


\concept{安全状态与不安全状态}
\begin{points}
  \item 安全状态指系统能找到一个安全序列，使所有进程都能按某种顺序获得所需资源并完成。
  \item 不安全状态不一定已经死锁，但可能发展为死锁。
  \item 死锁状态一定是不安全状态；不安全状态不一定是死锁状态。
\end{points}
\plain{死锁避免的核心不是“现在卡没卡死”，而是“这一步分配后系统还留不留得下安全退路”。}
\exam{常考安全、不安全、死锁三者关系，尤其“不安全不等于已死锁”。}


\concept{死锁、 安全状态、不安全状态的关系}
\begin{points}
  \item 安全状态一定不会死锁，因为至少存在一个安全序列。
  \item 不安全状态不等于已经死锁，但系统有可能进一步走向死锁。
  \item 死锁状态一定是不安全状态。
\end{points}
\plain{安全像“还有退路”，不安全像“退路没法保证了”，死锁则是“已经彻底堵死了”。}
\exam{常考判断题：不安全状态不一定死锁，这句话必须会。}


\concept{银行家算法}
\begin{points}
  \item 银行家算法用于死锁避免，假设系统知道每个进程的最大资源需求。
  \item 关键数据结构：Available 可用资源，Max 最大需求，Allocation 已分配，Need 尚需，且 $Need=Max-Allocation$。
  \item 资源请求满足 $Request_i\le Need_i$ 且 $Request_i\le Available$ 后，先试探性分配。
  \item 再执行安全性检查：寻找一个进程满足 $Need_i\le Work$，假设其完成并释放资源，更新 $Work=Work+Allocation_i$，直到所有进程完成或找不到下一个。
  \item 若试分配后仍安全，则真正分配；否则让进程等待。
\end{points}
\plain{银行家算法题的核心不是算表本身，而是证明“试分配之后还能排出一条让所有进程完成的安全序列”。}
\exam{常考 Available/Max/Allocation/Need 的含义、试分配步骤，以及安全序列的完整书写。}


\concept{银行家算法做题顺序}
\begin{points}
  \item 先检查 $Request_i\le Need_i$，否则请求本身非法。
  \item 再检查 $Request_i\le Available$，否则资源不足，进程等待。
  \item 只有前两步都通过，才允许试分配并修改 Available、Allocation、Need。
  \item 最后做安全性检查，若找不到安全序列，则撤销试分配。
\end{points}
\plain{老师在录音里反复强调的就是：别一上来就算安全序列，先做两道合法性检查。}
\exam{大题最容易扣分的点就是漏写前两步检查，或只写“安全”却不写安全序列。}


\concept{银行家算法与死锁检测的区别}
\begin{points}
  \item 银行家算法发生在资源分配之前，目标是避免系统进入不安全状态。
  \item 死锁检测发生在资源已经按通常方式分配之后，目标是判断当前是否已经发生死锁。
  \item 前者看到“不安全”就拒绝这次分配；后者若检测到死锁，还要继续指出死锁进程并考虑解除方案。
\end{points}
\plain{看题先分清是在问“这次能不能批”还是“现在是不是已经卡死”，这一步错了后面全会答偏。}
\exam{老师录音特别提醒要区分银行家算法和死锁检测，题型边界要先判断清楚。}


\concept{资源分配图}
\begin{points}
  \item 资源分配图用圆表示进程，用方框表示资源类型，方框内点表示资源实例。
  \item 请求边：进程 $\rightarrow$ 资源；分配边：资源实例 $\rightarrow$ 进程。
  \item 若每类资源只有一个实例，图中有环等价于发生死锁。
  \item 若资源类型有多个实例，有环只是死锁的必要条件，不一定死锁。
\end{points}
\plain{资源分配图题最易错的一句就是“有环是否一定死锁”，关键要先看资源类型是不是单实例。}
\exam{常考资源分配图中环的判定条件，尤其单实例和多实例资源时结论不同。}


\concept{死锁检测题的书写要求}
\begin{points}
  \item 若题目考死锁检测，要分清它不是银行家算法，不要求“避免进入不安全状态”。
  \item 检测出存在死锁时，答案不能只写“有死锁”，还应指出哪些进程陷入死锁。
  \item 当资源种类和实例较多时，可直接按检测过程判断，不一定非画资源分配图。
\end{points}
\plain{检测题关注“现在是不是已经卡死了”，而不是“这次要不要批资源”。}
\exam{老师录音明确强调：若存在死锁，一定写出发生死锁的进程集合。}


\concept{死锁检测的算法}
\begin{points}
  \item 检测算法与银行家算法中的安全性检查相似，但它处理的是当前已分配后的系统状态，不涉及 \code{Max} 和 \code{Need}，而是看当前 \code{Request} 能否被满足。
  \item 常用数据结构包括：\term{Available} 表示当前可用资源向量，\term{Allocation} 表示各进程已占资源矩阵，\term{Request} 表示各进程当前尚在请求的资源矩阵。
  \item 初始化时令 \code{Work = Available}；若某进程满足 $Allocation_i=0$，可先令 \code{Finish[i]=true}，否则 \code{Finish[i]=false}。
  \item 反复寻找满足 \code{Finish[i]=false} 且 $Request_i \le Work$ 的进程，若找到，就假设其完成并释放资源：\code{Work = Work + Allocation\_i}，再置 \code{Finish[i]=true}。
  \item 若最终仍存在 \code{Finish[i]=false} 的进程，则这些进程无法被当前系统资源链条推进，可判定它们陷入死锁。
\end{points}
\plain{检测算法的直觉是：不断试着找“现在就能被喂饱并完成”的进程，把它完成后释放出来的资源再拿去喂别人；最后一直喂不动的那批，就是被卡死的。}
\exam{常考它和银行家算法的区别：银行家算法看“试分配后安不安全”，检测算法看“在当前分配结果下还有没有进程能继续推进”。}


\concept{死锁检测与解除}
\begin{points}
  \item 死锁检测允许资源正常分配，周期性检查是否存在死锁。
  \item 检测后可通过撤销进程、回滚进程、剥夺资源等方式解除死锁。
  \item 选择牺牲进程时可考虑优先级、已运行时间、已占资源、还需资源、回滚代价等。
  \item 缺点是检测和恢复都有代价，且可能造成进程饥饿。
\end{points}
\plain{检测题要答到“查出来以后怎么办”，否则就只答了半题。}
\exam{常考撤销进程、资源剥夺、回滚恢复三类解除思路及其代价。}


\concept{综合方法}
\begin{points}
  \item 实际系统常把预防、避免、检测和人工干预结合使用，而不是全盘只用一种策略。
  \item 对关键且代价高的资源可用较严格策略，对一般资源可放宽限制，提高整体利用率。
\end{points}
\plain{现实系统不是教材单选题，往往是“重要资源严管，普通资源宽管”。}
\exam{常考综合方法的思想：在安全和效率之间折中。}
